\documentclass[11pt, a4paper]{article}

\usepackage[a4paper, top=2.2cm, bottom=2.2cm, left=2.4cm, right=2.4cm]{geometry}
\usepackage{fontenc}
\usepackage{inputenc}
\usepackage{microtype}
\usepackage{parskip}
\usepackage{booktabs}
\usepackage{tabularx}
\usepackage{array}
\usepackage{enumitem}
\usepackage{titlesec}
\usepackage{fancyhdr}
\usepackage{hyperref}
\usepackage{mdframed}
\usepackage{setspace}

\hypersetup{
  colorlinks=false,
  hidelinks
}

% Typography — serif font for direct Overleaf compilation
\renewcommand{\familydefault}{\rmdefault}
\setstretch{1.15}

% Section formatting
\titleformat{\section}{\large\bfseries}{}{0em}{}[\vspace{-4pt}\rule{\linewidth}{0.4pt}\vspace{2pt}]
\titleformat{\subsection}{\normalsize\bfseries}{}{0em}{}
\titleformat{\subsubsection}{\normalsize\itshape}{}{0em}{}
\titlespacing{\section}{0pt}{18pt}{6pt}
\titlespacing{\subsection}{0pt}{10pt}{3pt}
\titlespacing{\subsubsection}{0pt}{6pt}{2pt}

% List settings
\setlist[itemize]{noitemsep, topsep=3pt, leftmargin=1.4em}
\setlist[enumerate]{noitemsep, topsep=3pt, leftmargin=1.4em}

% Table column types
\newcolumntype{L}[1]{>{\raggedright\arraybackslash}p{#1}}
\newcolumntype{R}[1]{>{\raggedleft\arraybackslash}p{#1}}

% Header / footer
\pagestyle{fancy}
\fancyhf{}
\renewcommand{\headrulewidth}{0.3pt}
\fancyhead[L]{\small\textbf{MedMemory}: Product Requirements Document}
\fancyhead[R]{\small K V Naresh Karthigeyan}
\fancyfoot[C]{\small\thepage}

% Warning box
\newmdenv[
  leftmargin=0pt,
  rightmargin=0pt,
  innerleftmargin=10pt,
  innerrightmargin=10pt,
  innertopmargin=8pt,
  innerbottommargin=8pt,
  linewidth=0.6pt,
  frametitle={},
]{warnbox}

\begin{document}

% ─── TITLE BLOCK ─────────────────────────────────────────────────────────────
\begin{center}
  % {\large\textbf{MedMemory}}\\[4pt]
  {\large Product Requirements Document}\\[6pt]
  {\large\textbf{MedMemory}}\\[4pt]
  % \rule{\linewidth}{0.6pt}\\[4pt]
  {\small K V Naresh Karthigeyan}
\end{center}

\vspace{3pt}

% ─── 1. THE PROBLEM ──────────────────────────────────────────────────────────
\section{The Problem}

\textit{``My mother has three doctors. My father has two. And I live in a different city. The reports are in a plastic folder at home. Somebody photographs them sometimes.''}

\vspace{4pt}

The problem is not that families lack records, they do, but they are scattered, missing when needed the most, submitted to insurance claims, or uncategorized or outdated. It is that these records do not compound. Every new report lands in isolation. These can merely be a photograph in a WhatsApp chat, a PDF in someone's Downloads folder, a printout in a plastic sleeve. Nothing connects to what came before. The context lives only in the head of whoever last accompanied the patient.

\subsection{Who this breaks for}

\begin{itemize}
  \item \textbf{The adult child in another city.} Gets a panicked call at 11pm. No baseline. No medicine list. No idea what the cardiologist said three months ago.
  \item \textbf{The primary caregiver at home.} Carries all context mentally. Exhausting and fragile. One illness of their own and the system collapses.
  \item \textbf{The patient themselves.} Repeats their history to every new doctor from memory, which is unreliable and incomplete.
  \item \textbf{The consulting specialist.} Sees only what the patient brings. Usually, very little.
\end{itemize}

\subsection{The cost of fragmentation}

\vspace{4pt}
\begin{tabularx}{\linewidth}{@{}L{6.2cm} X@{}}
\toprule
\textbf{What happens} & \textbf{Concrete consequence} \\
\midrule
Reports do not connect across time & Lab Report values trend invisible across three reports from different labs \\
Medicine list lives in no single place & Cardiologist prescribes something conflicting with the patient's existing drug \\
Caregiver is the only source of truth & They travel for three days. Care degrades. \\
Emergency context unavailable & ER doctor asks blood group. Family member panics and guesses, or doesn't know, which wastes valuable time. \\
\bottomrule
\end{tabularx}

% ─── 2. SOLUTIONS CONSIDERED ─────────────────────────────────────────────────
\section{Solutions Considered}

\textbf{Option 1: A Smart folder / Drive integration System}
To Auto-organise uploads into Google Drive folders by date and patient. Simple. Fast to build. Does not solve the insight problem. You still cannot answer ``what changed this month'' or ``what medicines does she take.'' It is a tidier mess.

\textbf{Option 2: A Chatbot over medical records.}
Upload reports, ask questions in natural language. The problem: chatbots need a reason to exist. The actual need is ambient awareness, not Q\&A. A chatbot puts the cognitive load on the user to know what to ask. The caregiver calling from a different city does not have time to formulate queries.

\textbf{Option 3: Living medical memory that updates with each report. \textit{(Chosen.)}}
Each new report does not just get stored, but ratherr, it gets understood and integrated into a single patient record. The AI extracts structured entities (diagnoses, medicines, lab values, follow-ups), surfaces changes from prior state, flags conflicts, and generates a ready-to-use doctor brief. This is the only option that moves the cognitive load off the human. The output is not a file. It is a current understanding of the patient.

\subsection{What I deliberately chose not to not build:}

A generic patient portal or a pateint CRM for medicines (it adds login friction for caregivers who already face too much friction, and users are tired of sign ups that they forget soon), A symptom checker (medico-legal exposure and not the actual gap identified), A Native OCR from scratch (use an API and leading opensource models, they are surprisingly, computationally cheap, and the value is in the structuring layer, not text extraction).

% ─── 3. WHAT WE'RE BUILDING ──────────────────────────────────────────────────
\section{What we're Building}

MedMemory is a family health memory layer. You drop in a report. The AI reads it, extracts structured entities, and integrates them into a living patient record, a timeline that grows more useful with each upload rather than more cluttered.

\subsection{How someone uses it}

\vspace{4pt}
\begin{tabularx}{\linewidth}{@{}L{0.4cm} L{2.6cm} X@{}}
\toprule
\textbf{\#} & \textbf{Step} & \textbf{What happens} \\
\midrule
1 & Upload & Drop a PDF or photograph of a report. Works via WhatsApp forward too. \\
2 & AI processes & Extracts: diagnoses, medicines, lab values, doctor, date, follow-up instructions. Each entity confidence-scored. \\
3 & Review & Low-confidence extractions flagged for human verification. One tap to confirm or correct. \\
4 & Record updates & Timeline grows. Medicine list updates. Changes from prior reports surfaced. \\
5 & Act on it & Doctor brief ready to print. Emergency page always current. Shareable with family. \\
\bottomrule
\end{tabularx}

\vspace{6pt}
\subsection{Core outputs}

\vspace{4pt}
\begin{tabularx}{\linewidth}{@{}L{3.0cm} X@{}}
\toprule
\textbf{Output} & \textbf{Description} \\
\midrule
Living timeline & Chronological patient history. Each event links to its source report. \\
Unified medicine list & Across all doctors and hospitals. With confidence scores and conflict warnings. \\
Doctor brief & Printable one-pager: conditions, medicines, recent changes, questions to ask. \\
Emergency view & Blood group, allergies, medicines, contacts, insurance. Readable in ten seconds. \\
WhatsApp ingest & Future: parent forwards a report photo on WhatsApp. It processes and updates the record. \\
\bottomrule
\end{tabularx}

% ─── 4. WHERE AI IS LOAD-BEARING ─────────────────────────────────────────────
\section{Where the AI Is Load-Bearing}

The AI is not a chatbot bolted on. It is the structuring layer. Without it, the product is a file folder. With it, uploads become structured medical history.

\begin{itemize}
  \item \textbf{OCR and text extraction.} Turns photograph or PDF into readable text. Handles handwritten prescriptions and lab report formats across dozens of labs.
  \item \textbf{Medical entity extraction.} Finds and labels diagnoses, drug names and dosage, lab values with units and reference ranges, procedure dates, doctor names, hospital names, and follow-up instructions.
  \item \textbf{Conflict detection.} Flags when a new medicine conflicts with an existing one, or when a lab value diverges sharply from the patient's prior baseline.
  \item \textbf{Doctor brief generation.} Synthesises current state into a structured printable one-pager, including auto-generated questions based on recent lab changes.
  \item \textbf{Confidence scoring.} Assigns per-entity confidence. Anything below threshold is surfaced for human review. Not hidden.
\end{itemize}

% ─── 5. WHERE AI GETS IT WRONG ───────────────────────────────────────────────
\section{Where AI Will Get It Wrong}

% \begin{warnbox}
% \textbf{This section is not a disclaimer. It is a design constraint.} Every failure mode below has a corresponding product decision. AI errors are surfaced to users, not hidden. Human correction is the path of least resistance.
% % \end{warnbox}

% \vspace{6pt}
\begin{tabularx}{\linewidth}{@{}L{4.4cm} L{2.4cm} X@{}}
\toprule
\textbf{Failure mode} & \textbf{Likelihood} & \textbf{Product response} \\
\midrule
Handwritten prescription misread --- drug name garbled or dosage wrong &
High &
Low confidence triggers mandatory verification flag. Medicine not added to active list until confirmed. \\
\addlinespace
Lab value extracted without unit (12.3 --- but mg/dL or mmol/L?) &
Medium &
Unit is a required field in the entity model. Missing unit renders extraction incomplete, not silently used in trend. \\
\addlinespace
Medicine duplication --- same drug, different brand name across reports &
Medium &
Fuzzy match plus pharmacopoeia lookup. Probable duplicates flagged for user confirmation, not auto-merged. \\
\addlinespace
Hallucinated entity --- diagnosis extracted that is not in the report &
Low for structured reports; higher for narrative discharge summaries &
Confidence below 80\% flagged. All AI extractions shown with source-text highlight for user verification. \\
\addlinespace
Outdated medicine carried forward --- stopped drug still in active list &
High &
Medicines older than 90 days without reconfirmation shown as ``check status'' rather than active. \\
\addlinespace
OCR fails on low-resolution photo &
Medium &
Quality check prompts the user to re-upload or re-photograph. \\
\bottomrule
\end{tabularx}

% ─── 6. TEST CASES ───────────────────────────────────────────────────────────
\section{Test Cases and How to Test Them}

\subsection{Happy path (must pass before ship)}

\begin{enumerate}
  \item \textbf{Upload a clear lab report PDF.} Key values, units, date, and timeline event are correct.
  \item \textbf{Upload reports for two family members.} Entities go to the correct patient with no cross-contamination.
  \item \textbf{Generate a doctor brief after three reports.} Current state and lab trends are reflected, not only the last report.
\end{enumerate}

\subsection{Expected-failure cases (must handle gracefully)}

\begin{enumerate}
  \item \textbf{Upload a blurry prescription photo.} Show a quality warning and request a new upload. Do not add uncertain data silently.
  \item \textbf{Upload an unrecognised drug or duplicate report.} Flag it for verification and prevent duplicate entities.
  \item \textbf{Generate a brief with one report.} Generate it with a clear note that history is limited.
\end{enumerate}

\subsection{How to run this before ship}

Seed a test suite with twelve anonymised reports covering labs, discharge summaries, prescriptions, imaging, and degraded photographs. Score entity recall, precision, and confidence calibration. Manually audit every medicine extraction before launch.

\vspace{10pt}
\rule{\linewidth}{0.6pt}

\end{document}

\textit{Half a page, as requested. Sanity check, not a business plan.}
